\documentclass[11pt]{article}
\usepackage{xcolor}
\usepackage[italian]{babel}
\usepackage{amsmath, amsthm, amsfonts}
\usepackage[left=1.5cm, right=1.5cm, bottom=1cm, top=1cm]{geometry}

\usepackage{mdframed}
\mdfsetup{skipabove=1em,skipbelow=0em}
\theoremstyle{definition}
\newmdtheoremenv[nobreak=true]{definizione}{Definizione}
\newmdtheoremenv[nobreak=true]{proposizione}{Proposizione}
\newmdtheoremenv[nobreak=true]{lemma}{Lemma}
\newmdtheoremenv[nobreak=true]{teorema}{Teorema}
% \newmdtheoremenv[nobreak=true]{dimostrazione}{Dimostrazione}
\newtheorem{esempio}{Esempio}
\newtheorem{esercizio}{Esercizio}
\newtheorem{corollario}{Corollario}
\newtheorem*{corollario*}{Corollario}

\newtheorem*{notazione}{Notazione}
\surroundwithmdframed{notazione}
\newtheorem*{definizione*}{Definizione}
\surroundwithmdframed{definizione*}

\newtheorem*{esempio*}{Esempio}
\newtheorem*{proprieta}{Proprietà}

\theoremstyle{remark}
\newtheorem*{osservazione}{Osservazione}
\newtheorem*{dimostrazione}{Dimostrazione}

\newcommand{\Span}[1]{\left\langle{#1}\right\rangle}
\newcommand{\im}[1]{\text{Im}#1}

\pagestyle{empty}
\twocolumn

\begin{document}

\section*{Limiti notevoli}
\subsubsection*{$x\to 0$}
\begin{itemize}
    \item $\lim_{x\to0} \frac{\sin(x)}{x}=1$
    \item $\lim_{x\to0} \frac{\arcsin(x)}{x}=1$
    \item $\lim_{x\to0} \frac{\tan(x)}{x}=1$
    \item $\lim_{x\to0} \frac{\arctan(x)}{x}=1$
    \item $\lim_{x\to0} \frac{\cos(x)}{x^2}= \frac{1}{2}$
    \item $\lim_{x\to0} \frac{e^x-1}{x}=1$
    \item $\lim_{x\to0} \frac{\ln(1+x)}{x}=1$
\end{itemize}

\subsubsection*{$x\to \infty$}
\begin{itemize}
    \item $\lim_{x\to+\infty}(1+ \frac{\alpha}{x})^x=e^\alpha$
    \item $\lim_{x\to+\infty}{x^{ \frac{1}{x}}}=1$
    \item $\lim_{x\to+\infty}P(x)^{ \frac{1}{x}}=1$
\end{itemize}

\subsection*{Gerarchia infiniti}
\begin{equation*}
    (\log_Bn)^\alpha < n^\alpha < A^n < n! < n^n
\end{equation*}

\section*{Derivate}
\begin{itemize}
    \item $[ x^{\alpha}]'=\alpha x^{\alpha - 1}$
    \item $[ e^x ]'= e^x$
    \item $[ \sin(x)]'=\cos(x)$
    \item $[ \cos(x)]'=-\sin(x)$
    \item $[ \tan(x)]'= \frac{1}{\cos^2(x)}$
    \item $[ \arcsin(x)]'= \frac{1}{\sqrt{1-x^2}}$
    \item $[ \arccos(x)]'= -\frac{1}{\sqrt{1-x^2}}$
    \item $[ \arctan(x)]'=\frac{1}{1+x^2}$
\end{itemize}

\subsection*{Regole di derivazione}
\begin{itemize}
    \item $[kf(x)]'=kf'(x)$
    \item $[f(x)\pm g(x)]'=f'(x)\pm g'(x)$
    \item $[f(x)g(x)]'=f'(x)g(x) + f(x)g'(x)$
    \item $[\frac{f(x)}{g(x)}]'=\frac{f'(x)g(x) - f(x)g'(x)}{g(x)^2}$
    \item $[f(g(x))]'=f'(g(x))g'(x)$
\end{itemize}

\section*{Integrali}
\begin{itemize}
    \item $\int_{a}^{b} 1 dx= \left[x\right]_a^b$
    \item $\int_{a}^{b} x^n dx= \left[ \frac{x^{n+1}}{n+1}\right]_a^b$
    \item $\int_{a}^{b} \sin x dx= \left[-\cos x\right]_a^b$
    \item $\int_{a}^{b} \cos x dx= \left[\sin x\right]_a^b$
    \item $\int_{a}^{b} \frac{1}{\cos^2 x} dx= \left[\tan x\right]_a^b$
    \item $\int_{a}^{b} -\frac{1}{\sin^2 x} dx= \left[\cot x\right]_a^b$
    \item $\int_{a}^{b} \frac{1}{\sin^2 x} dx= \left[- \frac{1}{\tan x}\right]_a^b$
    \item $\int_{a}^{b} \frac{1}{1+x^2} dx= \left[\arctan x\right]_a^b$
    \item $\int_{a}^{b} \frac{1}{\sqrt{1-x^2}} dx= \left[\arcsin x\right]_a^b$
    \item $\int_{a}^{b} \frac{1}{\sqrt{1-x^2}} dx= \left[-\arccos x\right]_a^b$
    \item $\int_{a}^{b} \frac{1}{x^2} dx= \left[-\frac{1}{x}\right]_a^b$
    \item $\int_{a}^{b} \frac{1}{x} dx= \left[\ln|x|\right]_a^b$
    \item $\int_{a}^{b} e^x dx=\left[ e^x \right]_a^b$
\end{itemize}

\subsection*{Regole di integrazione}
\begin{itemize}
    \item $\int_{a}^{b} k\cdot f(x) dx = k\int_{a}^{b} f(x) dx$
    \item $\int_{a}^{b} f(x)\pm g(x) dx =\\= \int_{a}^{b} f(x) dx \pm \int_{a}^{b} g(x) dx$
\end{itemize}

\subsection*{Media integrale}
\begin{equation*}
    \Span{f} = \frac{1}{b-a}\int_{a}^{b} f(x) dx
\end{equation*}

\section*{Serie numeriche}
\subsection*{Serie geometrica}
\begin{equation*}
    \sum_{n=p}^{+\infty}x^n = x^p\sum_{i=0}^{+\infty}x^i=\begin{cases}
        +\infty & \text{ se } x>1 \\
        x^p\frac{1}{1-x} & \text{ se } -1<x<1 \\
        \not\exists & \text{ se } x\leq 1
    \end{cases}
\end{equation*}

\subsection*{Serie armoniche}
\begin{equation*}
    \sum_{n=0}^{+\infty} \frac{1}{n^{\alpha}}:\begin{cases}
        \text{Diverge} & \text{ se }\alpha\leq 1 \\
        \text{Converge} & \text{ se } \alpha > 1
    \end{cases}
\end{equation*}

\subsection*{Serie telescopiche}
Sia $a_n$ una serie tale che $a_n = b_{n+1} - b_n$. Si ha
\begin{equation*}
    S_n = \sum_{k=1}^{n}a_k = b_{n+1}-b_1
\end{equation*}
Quindi
\begin{equation*}
    \lim_{n\to+\infty}S_n = \sum_{n=1}^{+\infty}a_n = \lim_{n\to+\infty}b_{n+1}-b_1
\end{equation*}

\newpage
\subsection*{Criteri}
\begin{itemize}
    \item Radice
        \begin{equation*}
            \sum_{n=0}^{+\infty}a_n: \begin{cases}
                \text{Diverge} & \lim_{n\to+\infty}\sqrt[n]{a_n} > 1\\
                \text{Converge} & \lim_{n\to+\infty}\sqrt[n]{a_n} < 1
            \end{cases}
        \end{equation*}
    \item Rapporto
        \begin{equation*}
            \sum_{n=0}^{+\infty}a_n: \begin{cases}
                \text{Diverge} & \lim_{n\to+\infty} \frac{a_{n+1}}{a_n} > 1\\
                \text{Converge} & \lim_{n\to+\infty} \frac{a_{n+1}}{a_n} < 1\\            \end{cases}
        \end{equation*}
    \item Leibnitz
        \begin{equation*}
            \text{Se }\sum_{n=0}^{+\infty}(-1)^na_n: \begin{cases}
                a_n\geq 0 \\
                a_n \text{ è monotona decrescente} \\
                \lim_{n\to+\infty}a_n=0
            \end{cases} \\
        \end{equation*}
        Allora converge
\end{itemize}

\section*{Serie di potenze}
\subsection*{Raggio di convergenza}
Sia $S(n)=\sum_{n=0}^{+\infty}a_n(x-x_0)^n$
\begin{itemize}
    \item Se $\lim_{n\to+\infty}\sqrt{|a_n|}=l$ o $\lim_{n\to+\infty}|\frac{a_{n+1}}{a_n}|=l$:
        $\begin{cases}
            \rho = \frac{1}{l} & \text{ se } l\in(0, +\infty) \\
            \rho = \infty & \text{ se } l=0 \\
            \rho = 0 & \text{ se } l=\infty
        \end{cases}$
\end{itemize}
Calcolato $\rho$ si ha che $S(n)$ converge in $(x_0 - \rho, x_0 + \rho)$

\subsection*{Derivata nel centro}
Sia $S(x)=\sum_{n=0}^{+\infty}a_n(x-x_0)^n$
\begin{equation*}
    S^{(k)}(x_0) = k!a_k
\end{equation*}

\section*{Taylor}
\begin{equation*}
    T_n(f(x);x_0)=\sum_{k=0}^{n}\frac{f^{(k)}(x_0)}{k!}(x-x_0)^k
\end{equation*}
\subsection*{Funzioni elementari}
\begin{itemize}
    \item $e^x$
        \begin{equation*}
            T_n(e^x)=\sum_{k=0}^{n} \frac{x^k}{k!}
        \end{equation*}
    \item Seno/coseno
        \begin{gather*}
            T_{2n+1}(\sin(x);0) = \sum_{k=0}^{n} \frac{(-1)^kx^{2k+1}}{(2k+1)!}\\
            T_{2n}(\cos(x);0) = \sum_{k=0}^{n} \frac{(-1)^kx^{2k}}{(2k)!}
        \end{gather*}
    \item $\frac{1}{1\pm x}$
        \begin{gather*}
            T_n(\tfrac{1}{1-x};0)=\sum_{k=0}^{n}x^k \\
            T_n(\tfrac{1}{1+x};0)=\sum_{k=0}^{n}(-1)^kx^k \\
            T_{2n}(\tfrac{1}{1+x^2};0)=\sum_{k=0}^{n}(-1)^kx^{2k}
        \end{gather*}
    \item Logaritmo
        \begin{equation*}
            T_n(\log(1+x);0)= \sum_{k=1}^{n}(-1)^k \frac{x^k}{k}
        \end{equation*}
    \item $\arctan$
        \begin{equation*}
            T_{2n+1}(\arctan(x);0)= \sum_{k=0}^{n} \frac{(-1)^kx^{2k+1}}{2k+1}
        \end{equation*}
\end{itemize}


\end{document}
